\chapter{微分方程}\label{ch:b2:diffeq}

本科一年级已解决了可用公式求解的线性方程。本章
补充公式无法给出的内容：\emph{柯西--利普希茨定理}
--- 即 $y' = f(t, y)$ 的存在唯一性 --- 用
巴拿赫不动点定理证明，正如在
\cref{ch:b2:metric}中所承诺的；随后讨论\emph{线性
微分方程组} $X' = A(t)X + B(t)$ 的完整理论，以矩阵指数与
朗斯基行列式作为计算工具。

\section{柯西--利普希茨定理}

\begin{theorem}[柯西--利普希茨，全局利普希茨形式]\label{thm:b2:diffeq:cauchylipschitz}
设 $I$ 为闭区间，$f \colon I \times \R^n \to \R^n$ 连续
且\emph{关于第二变量满足利普希茨条件}，关于第一变量
一致地：$\norm{f(t, y) - f(t, z)} \leq k\,\norm{y - z}$ 对所有 $t
\in I$。于是对每个 $(t_0, y_0) \in I \times \R^n$，柯西
问题\index{柯西--利普希茨定理}
\[
y' = f(t, y), \qquad y(t_0) = y_0
\]
有且仅有一个解 $y \colon I \to \R^n$，它属于 $C^1$ 类。
\end{theorem}

\begin{proof}
\emph{等价改写。}连续的 $y$ 为该问题的解，当且仅当它
满足积分方程
\[
y(t) = y_0 + \int_{t_0}^{t} f\bigl(s, y(s)\bigr)\,\dd s
=: T(y)(t)
\]
（微积分基本定理双向适用；积分方程的连续解
自动为 $C^1$ 类）。

\emph{换范后的压缩映射。}在巴拿赫空间 $E =
C(I, \R^n)$ 上赋予\emph{加权}范数
\[
N(y) = \sup_{t \in I}\; \eu^{-2k\abs{t - t_0}}\,\norm{y(t)} ,
\]
（与上确界范数等价：权函数在
闭区间 $I$ 上有上下界，故 $E$ 仍完备），对 $y, z \in
E$，不妨设 $t \geq t_0$，估计：
\[
\norm{T(y)(t) - T(z)(t)}
\leq \int_{t_0}^{t} k\,\norm{y(s) - z(s)}\,\dd s
\leq k\,N(y - z)\int_{t_0}^{t} \eu^{2k(s - t_0)}\dd s
\leq \frac{N(y-z)}{2}\,\eu^{2k(t - t_0)} .
\]
将两边乘以 $\eu^{-2k(t - t_0)}$ 并取上确界（$t
< t_0$ 的情形类似）：$N\bigl(T(y) - T(z)\bigr) \leq \frac12 N(y -
z)$，故 $T$ 为 $\frac12$-压缩映射，且 $(E, N)$ 完备。
由巴拿赫不动点定理（见\cref{thm:b2:metric:banach}）得
唯一不动点，即唯一解。
\end{proof}

\begin{remark}
若 $f$ 仅为 $C^1$ 类（局部利普希茨），则定理
在\emph{局部}意义下成立，解为定义在极大开区间上的极大解；
解可能在有限时间内爆破（取 $y' = y^2$，$y(0) = 1$，则 $y(t)
= \frac{1}{1-t}$，在 $t = 1$ 处爆破）。全局利普希茨假设
是保证整个闭区间上解存在的关键。值得牢记的两条推论：
利普希茨场所决定的常微分方程，其解曲线\emph{永不相交}；
且线性齐次方程的解若在某点取零，
则该解恒为零。
\end{remark}

\begin{example}[唯一性是定理：一个会“漏”的场]\label{ex:b2:diffeq:leaky}
考虑 $y' = 2\sqrt{\abs y}$，且 $y(0) = 0$。零
函数是它的解；另一个解为
\[
y(t) = \begin{cases} 0 & t \leq 0,\\ t^2 & t \geq 0,
\end{cases}
\]
该函数是 $C^1$ 类的（两段在接合处的导数都为 $0$，
并满足 $y'(t) = 2t = 2\sqrt{t^2}$（$t > 0$）---
事实上，推迟起始时刻便给出一个解，对\emph{每个}
起始时刻 $c \geq 0$ 都有一个解，故经过
同一初值的解有无穷多个。这与
\cref{thm:b2:diffeq:cauchylipschitz}并不矛盾：在 $y = 0$ 附近，
\[
\frac{\abs{2\sqrt y - 2\sqrt z}}{\abs{y - z}}
= \frac{2}{\sqrt y + \sqrt z} \longrightarrow +\infty ,
\]
场关于 $y$ 不满足利普希茨条件，该定理在此不给出结论。
末了点出其物理含义：重力作用下排空的水桶
倒放——从空的状态出发，人们
无法判断它何时开始注水；常微分方程的确定性
恰是利普希茨条件，而非自然规律。
\end{example}

\section{线性微分方程组}

\begin{theorem}[线性微分方程组的结构]\label{thm:b2:diffeq:linear}
设 $A \colon I \to \mathcal{M}_n(\R)$ 与 $B \colon I \to \R^n$ 均
在区间 $I$ 上连续。对每个 $(t_0, X_0)$，问题
\[
X' = A(t)X + B(t), \qquad X(t_0) = X_0
\]
有且仅有一个定义在\emph{整个} $I$ 上的解。齐次
方程组（$B = 0$）的解构成向量空间 $\mathcal{S}_H$，其
维数恰为 $n$，且求值映射 $X \mapsto X(t_0)$ 是一个
同构 $\mathcal{S}_H \to \R^n$；一般解 $=$
特解 $+$ 齐次解。
\end{theorem}

\begin{proof}
对任一满足 $J \subseteq I$ 且包含 $t_0$ 的闭区间，$f(t, X) =
A(t)X + B(t)$ 连续，且关于 $X$ 满足利普希茨条件，常数 $k
= \sup_J \vertiii{A(t)}$（有限：在闭区间上连续）：
在 $J$ 上可应用\cref{thm:b2:diffeq:cauchylipschitz}；令 $J$
穷尽 $I$，由唯一性将这些解拼接为 $I$ 上的一个解。
解集与求值映射的线性显然成立；
求值映射由存在性得满射，由唯一性得
单射，故为双射，从而 $\dim \mathcal{S}_H = n$。仿射结构的讨论与
本科一年级的论证完全相同。
\end{proof}

\begin{example}[求值同构的具体形式]\label{ex:b2:diffeq:evaliso}
对 $y'' + y = 0$，将其视为方程组 $X' =
\begin{pmatrix} 0 & 1\\ -1 & 0\end{pmatrix}X$，取 $X = (y,
y')$：由定理，解空间为一个平面，且
$X \mapsto X(0) = (y(0), y'(0))$ 是到
$\R^2$ 的同构。解 $\cos$ 与 $\sin$ 的取值分别为 $(1, 0)$
与 $(0, 1)$ --- $\R^2$ 的标准基 ---，因此它们构成
解空间的一组基，因而\emph{任一}解都是
\[
y(t) = y(0)\cos t + y'(0)\sin t ,
\]
其中系数可直接由初值读出，无需
求解线性方程组。小结：选取
以标准基为初值的基本解组
（此处为 $\cos, \sin$），恰好是在选取
$\eu^{tA}$ 的各列；求值同构说明为何初值
刻画轨线 --- 即线性方程“确定性动力学”的
几何内容。
\end{example}

\begin{definition}[朗斯基行列式]\label{def:b2:diffeq:wronskian}
对齐次方程组的解 $X_1, \dots, X_n$，其
\emph{朗斯基行列式}\index{朗斯基行列式}为 $W(t) = \det\bigl(X_1(t),
\dots, X_n(t)\bigr)$。由上述同构，$W$ 要么
恒为零（族线性相关），要么处处不为零（
即\emph{基本解组}\index{基本解组}）；定量地，$W' = \operatorname{tr}\bigl(A(t)\bigr)
W$，于是
\[
W(t) = W(t_0)\,\exp\Bigl(\int_{t_0}^{t}
\operatorname{tr} A(s)\,\dd s\Bigr)
\quad \text{（刘维尔公式）}.
\]
\end{definition}

\begin{example}[在欧拉方程上验证刘维尔公式]\label{ex:b2:diffeq:eulerwronskian}
在 $\intoo{0}{\infty}$ 上，方程 $t^2y'' + ty' - y = 0$ 有
解 $y_1(t) = t$ 与 $y_2(t) = \frac1t$（代入即可验证）。
其朗斯基行列式为：
\[
W(t) = \det\begin{pmatrix} t & \tfrac1t\\[2pt]
1 & -\tfrac{1}{t^2}\end{pmatrix}
= -\frac1t - \frac1t = -\frac2t ,
\]
处处不为零，故为基本解组。现验证刘维尔公式：将方程化为
标准形式 $y'' + \frac1t\,y' - \frac{1}{t^2}\,y = 0$，其
友矩阵 $A(t) = \begin{pmatrix} 0 & 1\\ \frac{1}{t^2}
& -\frac1t\end{pmatrix}$ 的迹为 $-\frac1t$，于是
\[
W(t) = W(1)\exp\Bigl(-\int_1^t\frac{\dd s}{s}\Bigr)
= -2\,\eu^{-\ln t} = -\frac2t . \checkmark
\]
小结：在尚未知任何解时，刘维尔公式便预言了
朗斯基行列式的\emph{形状} --- 此处即 $W$ 必为
$\frac{c}{t}$；这正是
降阶法的依据
（见\cref{prop:b2:diffeq:secondorder}），由已知 $y_1$ 及
朗斯基行列式的形式，经一次求积即可定出 $y_2$。
\end{example}

\begin{proof}[刘维尔公式的证明]
$W(t) = \det M(t)$，其中 $M' = AM$。对行列式
视为各列的多重线性函数求导，
\[
W'(t) = \sum_j \det(X_1, \dots, X_j', \dots, X_n)
= \sum_j \det(X_1, \dots, AX_j, \dots, X_n) .
\]
现考虑映射 $(C_1, \dots, C_n) \mapsto \sum_j \det(C_1, \dots,
AC_j, \dots, C_n)$ 为 $n$ 重线性且交错的（若两列相等，
即 $C_i = C_k$，则 $j \notin \{i, k\}$ 的项均为零，
而 $j = i$ 与 $j = k$ 的两项经一次列交换后两两相消）：
由唯一性定理，
（\cref{thm:b2:linalg:detspace}）它等于 $c \cdot \det$，$c$
由标准列读出：$c = \sum_j \det(e_1, \dots, Ae_j, \dots,
e_n) = \sum_j a_{jj} = \operatorname{tr} A$。于是 $W' =
\operatorname{tr}\bigl(A(t)\bigr)W$，即标量线性
常微分方程，由本科一年级的公式求解。
\end{proof}

\section{常系数情形：矩阵指数}

\begin{theorem}\label{thm:b2:diffeq:matrixexp}
对 $A \in \mathcal{M}_n(\R)$（或 $\C$），指数 $\eu^{tA}
= \sum_k \frac{(tA)^k}{k!}$（见\cref{ex:b2:nvs:matrixexp}）满足：
$t \mapsto \eu^{tA}$ 是 $C^1$ 的（事实上是 $C^\infty$ 的），
且有\index{矩阵指数}
\[
\frac{\dd}{\dd t}\,\eu^{tA} = A\,\eu^{tA} = \eu^{tA}A ,
\qquad
\eu^{(s+t)A} = \eu^{sA}\,\eu^{tA},
\qquad
(\eu^{A})^{-1} = \eu^{-A} ;
\]
且 $\eu^{A + B} = \eu^A\eu^B$，\emph{当} $AB = BA$ 时成立。柯西
问题 $X' = AX$，$X(0) = X_0$ 的唯一解为 $X(t) =
\eu^{tA}X_0$；若含源项，则\emph{常数变易}
公式成立：
\[
X(t) = \eu^{(t - t_0)A}X_0 + \int_{t_0}^{t} \eu^{(t-s)A}B(s)\,\dd
s .
\]
\end{theorem}

\begin{proof}
\emph{可微性：}级数 $\sum \frac{t^kA^k}{k!}$ 与
其逐项求导所得级数 $\sum \frac{t^{k-1}A^k}{(k-1)!} = A\sum
\frac{(tA)^{k-1}}{(k-1)!}$ 在任一闭区间上正规收敛
（范数 $\leq \frac{(\abs t\,\vertiii A)^k}{k!}$），可逐项
求导（见\cref{thm:b2:funcseq:seriestransfer}，向量值情形）。
两种次序 $A\eu^{tA}$ 与 $\eu^{tA}A$ 相等，因为每个
部分和均与 $A$ 可交换。

\emph{群律：}对于可交换的 $A, B$，两个指数级数的柯西乘积
依二项式定理重新整理，与
\cref{ex:b2:series:exp}所示完全相同（在巴拿赫
代数中，绝对收敛性保证了这一步）：$\eu^{A+B} = \eu^A\eu^B$；取 $B = sA$，
便得单参数群律，取 $B = -A$ 则得逆元。

\emph{柯西问题：}$X(t) = \eu^{tA}X_0$ 是其解（求导即可验证）；
唯一性由\cref{thm:b2:diffeq:linear}可得。常数变易法：
设 $Y(t) = \eu^{-tA}X(t)$；求导得 $Y' = \eu^{-tA}(X' -
AX) = \eu^{-tA}B(t)$；从 $t_0$ 到 $t$ 积分，
再左乘 $\eu^{tA}$ 即可。
\end{proof}

\begin{method}[计算 $\eu^{tA}$]\label{met:b2:diffeq:computeexp}
化约 $A$（见\cref{ch:b2:reduction}）：若 $A = PDP^{-1}$ 为对角形，
则 $\eu^{tA} = P\,\eu^{tD}P^{-1}$，其中 $\eu^{tD}$ 为对角矩阵，
其对角元为 $\eu^{t\lambda_i}$；一般情形用邓福德分解 $A = D + N$
（二者可交换）：$\eu^{tA} = \eu^{tD}\,\eu^{tN}$，其中 $\eu^{tN}$ 是
关于 $t$ 的\emph{多项式}（幂零性使级数截断）。复
特征值成对出现，对应“旋转乘指数”型的块
（见\cref{exo:b2:diffeq:5}）。
\end{method}

\begin{remark}[常见陷阱]
\emph{(i) 不可交换时 $\eu^{A+B} \neq \eu^A\eu^B$：}
取 $A = \begin{pmatrix} 0 & 1\\ 0 & 0\end{pmatrix}$，$B =
\begin{pmatrix} 0 & 0\\ 1 & 0\end{pmatrix}$。则 $\eu^A = I +
A$，$\eu^B = I + B$（由幂零性），于是
\[
\eu^A\eu^B = \begin{pmatrix} 2 & 1\\ 1 & 1\end{pmatrix},
\qquad\text{而}\qquad
\eu^{A+B} = \cosh(1)\,I + \sinh(1)\,(A + B)
= \begin{pmatrix} \cosh 1 & \sinh 1\\ \sinh 1 & \cosh 1
\end{pmatrix},
\]
利用 $(A+B)^2 = I$，且 $\cosh 1 \approx 1.54 \neq 2$。
\cref{thm:b2:diffeq:matrixexp}所述的群律确有
前提条件。\emph{(ii) 线性问题中的非线性直觉误区：}
具有连续系数的\emph{线性}系统的解
在整个区间上均有定义
（见\cref{thm:b2:diffeq:linear}） --- 若某候选解
在 $I$ 内爆破，则方程不是线性的，或计算有误；
反之，对于非线性方程，
不经论证，绝不可断言解整体存在（$y' = y^2$）。
\emph{(iii) 除以未知量：}对 $y' = y(1-y)$ 分离变量时，
会悄然丢掉常数解 $0$ 与
$1$ --- 恰好组织起相线的那些解
（见\cref{exo:b2:diffeq:3}）；应先列出常数解。
\emph{(iv) 初值确定的是向量，而非标量：}一个 $n$ 阶
标量方程需要 $n$ 个条件（$y, y', \dots$ 在
$t_0$ 处）；仅匹配 $y(t_0)$ 会留下一个
$(n-1)$ 参数族，是“丢失”
常数的经典来源。
\end{remark}

\begin{example}[用邓福德分解计算 $3\times3$ 指数]\label{ex:b2:diffeq:dunford3}
求解 $X' = AX$，其中 $A = \begin{pmatrix} 2 & 1 & 0\\ 0 & 2 &
0\\ 0 & 0 & 3\end{pmatrix}$。按块作邓福德分解：$A = D + N$
其中 $D = \operatorname{diag}(2, 2, 3)$，$N = E_{12}$，二者
可交换（$N$ 位于特征值 $2$ 的块内），且 $N^2 =
0$：
\[
\eu^{tA} = \eu^{tD}\,\eu^{tN}
= \begin{pmatrix}
\eu^{2t} & t\,\eu^{2t} & 0\\
0 & \eu^{2t} & 0\\
0 & 0 & \eu^{3t}
\end{pmatrix} .
\]
逐列读出通解：$X(t) =
\bigl(\eu^{2t}(x_0 + ty_0),\ \eu^{2t}y_0,\
\eu^{3t}z_0\bigr)$。验算：当 $t = 0$ 时矩阵为
$I$；其行列式为 $\eu^{7t} =
\eu^{t\operatorname{tr}A}$，符合刘维尔公式；而
因子 $t$ 恰好出现在特征值 $2$
亏损的地方。要点：多项式乘指数函数
并非需要死记的猜测 --- 它们正是截断后的级数
$\eu^{tN}$，其次数以幂零指数为界，
绝不会更高。
\end{example}

\begin{method}[从头到尾求解 $X' = AX + B(t)$]\label{met:b2:diffeq:solvelinear}
\begin{enumerate}
  \item 求 $A$ 的谱，再用
        \cref{met:b2:diffeq:computeexp}求出 $\eu^{tA}$（对角化；或
        用邓福德分解，如\cref{ex:b2:diffeq:dunford3}；或
        利用多项式技巧，例如 $A^2 = -I$）。
  \item 特解：常数变易法
        $\int_{t_0}^t\eu^{(t-s)A}B(s)\dd s$ 总是可行；
        对于指数多项式型的 $B$，取形式相同的试探解
        （共振时次数升高，
        见\cref{exo:b2:diffeq:10}）更快。
  \item 通解 $= \eu^{(t-t_0)A}X_0 +$ 特解；
        初值\emph{最后}代入，
        在完整公式上进行。
  \item 验算：$X(t_0)$ 正确；齐次部分的
        增长与特征值实部相符
        （见\cref{exo:b2:diffeq:8}）；基解矩阵的 $\det$
        符合刘维尔公式。
\end{enumerate}
\end{method}

\begin{example}[相图示例]\label{ex:b2:diffeq:phaseportrait}
$X' = AX$，其中 $A = \begin{pmatrix} 0 & 1\\ -1 & 0\end{pmatrix}$：
$A^2 = -I$，于是级数分解为
\[
\eu^{tA} = (\cos t)\,I + (\sin t)\,A
= \begin{pmatrix} \cos t & \sin t\\ -\sin t & \cos t
\end{pmatrix} :
\]
轨道是沿顺时针方向运行的圆 --- 即
以一阶形式写出的调和振子 $x'' + x = 0$。特征值 $\pm\iu$ 位于
虚轴上，此为\emph{中心}。更一般地，
$A$ 的特征值的实部决定 $\norm{X(t)}$ 的增长或衰减
（见\cref{exo:b2:diffeq:8}）。
\end{example}

\begin{omfigure}
\begin{tikzpicture}
\begin{axis}[omaxis, width=6.4cm, height=6.4cm,
    xmin=-2.3, xmax=2.3, ymin=-2.3, ymax=2.3,
    xtick={-1,1}, ytick={-1,1}]
  \addplot[omDef, thick, domain=0:360, samples=90, smooth]
    ({cos(x)}, {-sin(x)});
  \addplot[omDef!60, thick, domain=0:360, samples=90, smooth]
    ({1.7*cos(x)}, {-1.7*sin(x)});
  \draw[-Stealth, omThm, thick] (axis cs:1,0) -- (axis cs:0.995,-0.12);
  \draw[-Stealth, omThm, thick] (axis cs:1.7,0) -- (axis cs:1.695,-0.12);
\end{axis}
\end{tikzpicture}
\hspace{6mm}
\begin{tikzpicture}
\begin{axis}[omaxis, width=6.4cm, height=6.4cm,
    xmin=-2.3, xmax=2.3, ymin=-2.3, ymax=2.3,
    xtick={-1,1}, ytick={-1,1}]
  % A = diag(-1,-2)-type stable node: trajectories y = c x^2
  \addplot[omMeth, thick, domain=-1.5:1.5, samples=60] {x^2};
  \addplot[omMeth, thick, domain=-1.5:1.5, samples=60] {-x^2};
  \addplot[omMeth, thick, domain=-2.2:2.2, samples=2] {0};
  \addplot[omMeth!60, thick, domain=-1.2:1.2, samples=60] {2*x^2};
  \draw[-Stealth, omThm, thick] (axis cs:1.2,1.44) --
    (axis cs:1.05,1.1);
  \draw[-Stealth, omThm, thick] (axis cs:-1.2,-1.44) --
    (axis cs:-1.05,-1.1);
\end{axis}
\end{tikzpicture}

{\small 两个线性相图。左：\emph{中心}
（特征值 $\pm\iu$） --- 调和振子的闭合圆轨道。
右：\emph{稳定结点}（特征值 $-1, -2$） ---
所有轨道沿慢
特征方向相切地落入原点。}
\end{omfigure}

\begin{omfigure}
\begin{tikzpicture}
\begin{axis}[omaxis, width=9.6cm, height=6.4cm,
    xmin=-3.2, xmax=3.2, ymin=-1.6, ymax=2.6,
    xtick={0}, ytick={0},
    xlabel={$\tau = \operatorname{tr}A$},
    ylabel={$\delta = \det A$}]
  \addplot[omThm, very thick, domain=-3:3, samples=100]
    {x^2/4};
  \node[font=\small, omThm] at (axis cs:2.45,2.1)
    {$\delta = \tau^2/4$};
  \node[font=\small] at (axis cs:0,-0.9) {鞍点};
  \node[font=\small] at (axis cs:-2.45,0.7)
    {稳定结点};
  \node[font=\small] at (axis cs:2.45,0.7)
    {不稳定结点};
  \node[font=\small] at (axis cs:-1.05,1.9)
    {稳定焦点};
  \node[font=\small] at (axis cs:1.05,1.9)
    {不稳定焦点};
  \node[font=\small, omDef] at (axis cs:0,2.35) {中心};
  \draw[omDef, very thick] (axis cs:0,0) -- (axis cs:0,2.55);
\end{axis}
\end{tikzpicture}

{\small $X' = AX$ 在 $2$ 维情形下的迹--行列式平面：
水平轴下方为鞍点；水平轴与抛物线 $\delta = \tau^2/4$ 之间为结点；
抛物线内部为焦点；
正 $\delta$ 轴上为中心。周末问题证明了这一分类，
并沿其中一条竖直线 --- 阻尼振子 ---
考察其进入共振的过程。}
\end{omfigure}

\begin{remark}[应用]
线性系统是处理一切非线性问题的局部模型：
在平衡点附近，光滑向量场（在双曲情形下）的行为
与其线性化的行为一致，其相图由迹--行列式
平面刻画。周末问题完整地讨论了振子的故事 ---
阻尼、受迫、共振，以及变系数情形下的
施图姆比较定理 ---
这正是减震器、交流电路与
谱隙背后共同的数学。本科三年级分册
将在流形上回到定性理论（流、稳定性、
首次积分）。
\end{remark}

\section{变系数的二阶方程}

\begin{proposition}\label{prop:b2:diffeq:secondorder}
方程 $y'' + a(t)y' + b(t)y = c(t)$（$a, b, c$ 在 $I$ 上连续）
等价于系统 $X' = A(t)X + B(t)$，其中 $X = (y, y')$：
在整个 $I$ 上解存在且唯一，初值 $(y(t_0),
y'(t_0))$ 可任意给定；齐次解构成一个平面。若已知一个
处处不为零的齐次解 $y_1$，则可通过\emph{降阶}求出与之独立的
第二个解：令 $y = y_1 z$，齐次方程化为关于 $z'$ 的
一阶方程，
再由求积分法解出。
\end{proposition}

\begin{proof}
由系统形式及\cref{thm:b2:diffeq:linear}
即可得到全部结构性结论。降阶：代入 $y = y_1z$，
\[
y_1 z'' + (2y_1' + a y_1)z' + \underbrace{(y_1'' + ay_1' +
by_1)}_{=\,0}\,z = 0 :
\]
这是关于 $u = z'$ 的一阶线性方程，可由本科一年级的
公式解出；对 $u$ 积分得 $z$，从而 $y_2 = y_1 z$，
它与 $y_1$ 独立，只要 $z$ 非常数。
\end{proof}

\begin{example}\label{ex:b2:diffeq:lowering}
$t^2y'' - 2y = 0$ 在 $\intoo{0}{\infty}$ 上：$y_1 = t^2$ 是一个
解。令 $y = t^2z$，由 $y' = t^2z' + 2tz$ 及 $y''
= t^2z'' + 4tz' + 2z$，得
\[
t^2y'' - 2y = t^4 z'' + 4t^3z' = 0,
\qquad\text{即}\qquad \frac{z''}{z'} = -\frac4t :
\]
$z' = t^{-4}$（相差一个常数），$z = -\frac{1}{3t^3}$，于是 $y_2
= t^2z = -\frac{1}{3t}$。通解为 $y = \alpha t^2 +
\frac{\beta}{t}$。
\end{example}

\section{练习}

\begin{exercise}[$\star$]\label{exo:b2:diffeq:1}
求解 $X' = AX$，$X(0) = (1, 0)^{\mathsf T}$，其中 $A =
\begin{pmatrix} 1 & 1\\ 0 & 2\end{pmatrix}$（对角化），以及 $A =
\begin{pmatrix} 2 & 1\\ 0 & 2 \end{pmatrix}$（用邓福德分解）。
\end{exercise}

\begin{exercise}[$\star$]\label{exo:b2:diffeq:2}
哪些柯西问题在 $\R$ 上有唯一的整体解，可由
\cref{thm:b2:diffeq:cauchylipschitz}判定？$y' = \sin(ty)$；$\;y' =
y^2$；$\;y' = \abs y$。对最后一个，分别在 $y(0) =
0$ 与 $y(0) = 1$ 下显式求解。
\end{exercise}

\begin{exercise}[$\star$]\label{exo:b2:diffeq:3}
证明：对 $y' = f(t,y)$ 的两个不同的极大解（$f$
关于 $y$ 满足利普希茨条件），在同一时刻永不取相同的值，
并由此推出：$y' = y(1 - y)$ 的解若始于
$\intoo{0}{1}$ 中，则永远留在 $\intoo{0}{1}$ 中。
\end{exercise}

\begin{exercise}[$\star\star$]\label{exo:b2:diffeq:4}
计算 $\eu^{tA}$，其中 $A = \begin{pmatrix} 3 & 1\\ -1 &
1\end{pmatrix}$ \emph{（用邓福德分解：$(A - 2I)^2 = 0$）}，并求解 $X' =
AX + \begin{pmatrix} \eu^{2t}\\ 0\end{pmatrix}$，$X(0) = 0$，用
常数变易法。
\end{exercise}

\begin{exercise}[$\star\star$]\label{exo:b2:diffeq:5}
设 $A = \begin{pmatrix} \alpha & -\beta\\ \beta &
\alpha\end{pmatrix}$，证明 $\eu^{tA} =
\eu^{\alpha t}\begin{pmatrix} \cos\beta t & -\sin\beta t\\
\sin\beta t & \cos\beta t\end{pmatrix}$ --- 螺线轨迹 ---
两种方法：利用级数（写 $A = \alpha I + \beta J$，$J^2 =
-I$），以及借助复数形式 $z' = (\alpha +
\iu\beta)z$。
\end{exercise}

\begin{exercise}[$\star\star$]\label{exo:b2:diffeq:6}
（格朗沃尔引理）设 $u$ 为非负连续函数，且 $u(t)
\leq C + k\int_{t_0}^{t} u(s)\,\dd s$ 在 $\intco{t_0}{T}$ 上成立。证明
$u(t) \leq C\,\eu^{k(t - t_0)}$ \emph{（对 $v(t) =
\eu^{-kt}\int_{t_0}^t u$ 求导）}。再由此推出
柯西--利普希茨定理中的唯一性，以及解的连续依赖性：$\norm{y(t) -
z(t)} \leq \norm{y_0 - z_0}\,\eu^{k\abs{t - t_0}}$，
对初值不同的两个解成立。
\end{exercise}

\begin{exercise}[$\star\star$]\label{exo:b2:diffeq:7}
已知 $y_1(t) = \frac{\sin t}{t}$ 是 $ty'' + 2y' + ty =
0$ 在 $\intoo{0}{\pi}$ 上的一个解，求与之线性无关的第二个解，
用降阶法，并给出通解。
\end{exercise}

\begin{exercise}[$\star\star\star$]\label{exo:b2:diffeq:8}
设 $A \in \mathcal{M}_n(\C)$，其全部特征值的
（严格）实部为负。证明 $X' = AX$ 的每个解都趋于
$0$（当 $t \to +\infty$ 时），且具有指数速率：$\norm{X(t)}
\leq C\,\eu^{-\alpha t}$，对某个 $\alpha > 0$。
\emph{（三角化；从最后一行向上逐行处理三角方程组，
或用邓福德分解：$\eu^{tA} = \eu^{tD}\eu^{tN}$，其中
$\norm{\eu^{tD}} \leq \eu^{-\alpha' t}$，且 $\eu^{tN}$ 为
关于 $t$ 的多项式。）}
\end{exercise}

\begin{exercise}[$\star\star\star$]\label{exo:b2:diffeq:9}
（线性增长时不会在有限时间内逃逸）设 $f$ 连续，
满足 $\norm{f(t, y)} \leq a\norm y + b$，在
$\intco{0}{\infty} \times \R^n$ 上关于 $y$ 局部利普希茨。利用
积分形式的格朗沃尔引理（见\cref{exo:b2:diffeq:6}），证明
极大解都是整体解（定义在
整个 $\intco{0}{\infty}$ 上）。
\end{exercise}

\begin{exercise}[$\star$]\label{exo:b2:diffeq:10}
求解 $y'' - 3y' + 2y = \eu^{t}$：先求齐次方程的解，再求
形如 $\alpha t\,\eu^{t}$ 的特解 \emph{（朴素猜测
$\alpha\eu^t$ 为何失效？）}；通解，以及
满足 $y(0) = y'(0) = 0$ 的解。
\end{exercise}

\begin{exercise}[$\star\star$]\label{exo:b2:diffeq:11}
对如下若尔当块计算 $\eu^{tA}$：
\[
A = \begin{pmatrix} \lambda & 1 & 0\\ 0 & \lambda & 1\\
0 & 0 & \lambda\end{pmatrix},
\]
并描述 $X' = AX$ 的全部解：指数函数乘以
次数至多为 $2$ 的多项式向量。这一多项式次数
从何而来？
\end{exercise}

\begin{exercise}[$\star\star\star$]\label{exo:b2:diffeq:12}
（周期强迫，周期响应）设 $A \in
\mathcal{M}_n(\R)$，$B \colon \R \to \R^n$ 连续且
以 $T$ 为周期。
\begin{enumerate}
  \item 证明 $X' = AX + B(t)$ 的解
        是 $T$ 周期的，当且仅当 $X(T) = X(0)$
        \emph{（比较 $X(\cdot + T)$ 与 $X$）}。
  \item 证明 $\eu^{TA}$ 的特征值恰为
        $\eu^{T\lambda}$, $\lambda \in \operatorname{Sp}A$
        \emph{（在 $\C$ 上三角化）}。由此可得：若
        $A$ 的特征值都不属于 $\frac{2\iu\pi}{T}\Z$，则
        $I - \eu^{TA}$ 可逆。
  \item 在该假设下，证明该方程组恰有一个
        $T$ 周期解，且
        \[
        X(0) = \bigl(I - \eu^{TA}\bigr)^{-1}
        \int_0^{T}\eu^{(T-s)A}B(s)\,\dd s .
        \]
        对于谐振子，被排除的情形
        对应什么？（周末问题的答案为：
        共振。）
\end{enumerate}
\end{exercise}

\section{问题：振动、共振与施图姆
比较定理}

\begin{problem}\label{pb:b2:diffeq:1}
一个方程支配着机械与电气世界：
\[
x'' + 2\zeta\omega\,x' + \omega^2 x = F(t),
\qquad \omega > 0,\ \zeta \geq 0 .
\]
本题将完整研究它 --- 借助平面线性系统的
迹--行列式分类、三种阻尼情形、
周期外力下的稳态响应，包括其
\emph{共振}\index{共振}峰值
与共振灾变 --- 随后撇开常系数情形，
转而讨论\emph{施图姆分离定理与比较
定理}\index{施图姆分离定理与比较定理}，它们控制
$y'' + q(t)y = 0$ 之解的零点，且完全不需要
任何显式公式。

\medskip
\noindent\textbf{第一部分 --- 迹--行列式平面。}设
$A \in \mathcal M_2(\R)$，$\tau = \operatorname{tr}A$，$\delta
= \det A$，$\Delta = \tau^2 - 4\delta$。
\begin{enumerate}
  \item 证明 $A$ 的特征值为
        $\frac{\tau\pm\sqrt\Delta}{2}$，并分类：两个实
        特征值符号相反，当且仅当 $\delta < 0$；实
        特征值同号，当且仅当 $\delta > 0$，$\Delta
        \geq 0$（与 $\tau$ 同号）；共轭复数对，当且仅当
        $\Delta < 0$（实部为 $\frac\tau2$）。
  \item （鞍点，$\delta < 0$）设特征值为 $\mu < 0 <
        \lambda$，特征向量为 $v_\pm$，写出通解，并
        描述轨线：两条稳定射线和
        两条不稳定射线，其余轨道均渐近于两组
        射线。为什么除 $0$ 外，没有任何解
        能在整个 $\R$ 上保持有界？
  \item （结点，$\delta > 0$，$\Delta > 0$）对于 $\mu < \lambda
        < 0$：证明每个非零解都趋于 $0$，并
        说明除快轴上的轨道外，其余轨道均
        \emph{与慢特征方向相切}地趋于原点
        \emph{（比较 $\eu^{\mu t}$ 与 $\eu^{\lambda t}$）}。
  \item （螺旋线与中心，$\Delta < 0$）记特征值为
        $\alpha \pm \iu\beta$，利用
        \cref{exo:b2:diffeq:5}（经实基变换后，
        在该一般情形下承认，或对第二部分中的系统加以证明，
        这些系统即下文所用者）来描述
        轨道：当 $\alpha =
        \frac\tau2 < 0$ 时螺旋线收敛，$\tau > 0$ 时
        螺旋线发散，$\tau = 0$ 时为闭曲线（中心）。
  \item （边界情形）对于二重特征值（$\Delta =
        0$）：证明 $\eu^{tA} = \eu^{\lambda t}(I + tN)$，其中 $N
        = A - \lambda I$ 幂零，并区分星形结点
        （$N = 0$）与非正常结点（$N \neq 0$）。
        在本章图形的迹--行列式图中概括
        第一部分的内容。
\end{enumerate}

\medskip
\noindent\textbf{第二部分 --- 阻尼振子。}此时 $F =
0$：$x'' + 2\zeta\omega x' + \omega^2x = 0$，即\ $X' = AX$
其中 $A = \begin{pmatrix} 0 & 1\\ -\omega^2 &
-2\zeta\omega\end{pmatrix}$。
\begin{enumerate}[resume]
  \item 计算 $\tau, \delta, \Delta$，并把三种
        阻尼情形标在迹--行列式平面上：
        \emph{欠阻尼} $0 < \zeta < 1$（稳定螺旋线），
        \emph{临界阻尼} $\zeta = 1$（二重特征值），
        \emph{过阻尼} $\zeta > 1$（稳定
        结点）；$\zeta = 0$ 对应中心。
  \item 显式求出三种情形的解：
        \[
        \zeta < 1:\ \eu^{-\zeta\omega t}\bigl(a\cos\omega_d t
        + b\sin\omega_dt\bigr),\ \omega_d =
        \omega\sqrt{1-\zeta^2};
        \qquad
        \zeta = 1:\ (a + bt)\,\eu^{-\omega t};
        \]
        $\zeta > 1$：两个实指数函数。定义
        伪周期 $\frac{2\pi}{\omega_d}$，并证明
        $\abs x$ 的相邻极大值之比为
        常数 $\eu^{-2\pi\zeta/\sqrt{1-\zeta^2}}$（即
        对数衰减率）。
  \item （闭门器原理）对于 $\zeta \geq 1$，
        衰减速率由最慢的特征值决定，记为
        $\lambda_{\mathrm{slow}} = -\omega\bigl(\zeta -
        \sqrt{\zeta^2-1}\bigr)$。证明
        $\abs{\lambda_{\mathrm{slow}}} =
        \frac{\omega}{\zeta + \sqrt{\zeta^2 - 1}}$ 是一个
        关于 $\zeta \geq 1$ 的\emph{递减}函数：
        临界阻尼 $\zeta = 1$ 给出最快的
        非振荡回到静止状态。
  \item （能量）设 $E(t) = \frac12x'^2 +
        \frac12\omega^2x^2$。证明 $E' = -2\zeta\omega\,x'^2
        \leq 0$，并推出当 $\zeta > 0$ 时方程
        没有非零周期解\emph{（若有周期，则 $E$
        为常数，从而 $x' \equiv 0$）}。
  \item 用两句话解释，为何中心 $\zeta = 0$
        是\emph{结构上脆弱}的：任何 $\zeta > 0$，
        无论多小，都会破坏周期性 --- 并指出这在迹--行列式平面上
        的表现（中心直线
        没有内点）。
\end{enumerate}

\medskip
\noindent\textbf{第三部分 --- 受迫振动与
共振。}此时 $F(t) = F\cos(\gamma t)$，其中 $F, \gamma > 0$。
\begin{enumerate}[resume]
  \item （$\zeta > 0$：稳态）试取 $x_p =
        \Re\bigl(z\,\eu^{\iu\gamma t}\bigr)$，并证明
        \[
        z = \frac{F}{\omega^2 - \gamma^2 +
        2\iu\zeta\omega\gamma},
        \qquad
        A(\gamma) := \abs z = \frac{F}{\sqrt{(\omega^2 -
        \gamma^2)^2 + 4\zeta^2\omega^2\gamma^2}} ,
        \]
        并写出 $x_p = A(\gamma)\cos(\gamma t - \varphi)$，
        其中 $\tan\varphi =
        \frac{2\zeta\omega\gamma}{\omega^2-\gamma^2}$。
  \item 证明：\emph{每一个}解都是 $x_p$ 加上一个
        第二部分中的瞬态项，它趋于 $0$；无论
        初值如何，系统都会锁定到
        稳态 --- 振幅为 $A(\gamma)$，相位滞后
        $\varphi$。
  \item （共振曲线）求 $A$ 的最大值：证明
        $A(\gamma)$ 有内部极大值，当且仅当 $\zeta <
        \frac{1}{\sqrt2}$ 时，极大值点与极大值分别为
        \[
        \gamma_* = \omega\sqrt{1 - 2\zeta^2},
        \qquad
        A(\gamma_*) =
        \frac{F}{2\zeta\omega^2\sqrt{1-\zeta^2}} ,
        \]
        并且当 $\zeta$ 很小时，峰值使静态响应
        $A(0) = \frac F{\omega^2}$ 放大到
        $\approx \frac{1}{2\zeta}$ 倍。
  \item （$\zeta = 0$，非共振）对于 $\gamma \neq \omega$，
        证明满足 $x(0) = x'(0) = 0$ 的解为
        \[
        x(t) = \frac{F}{\omega^2 -
        \gamma^2}\bigl(\cos\gamma t - \cos\omega t\bigr)
        = \frac{2F}{\omega^2-\gamma^2}
        \sin\frac{(\omega-\gamma)t}{2}
        \sin\frac{(\omega+\gamma)t}{2} :
        \]
        有界，且出现\emph{拍}现象 --- 即
        一个快振荡位于慢包络之下 --- 当 $\gamma$ 接近
        $\omega$ 时出现。
  \item （$\zeta = 0$，共振）对于 $\gamma = \omega$，证明
        $x_p(t) = \frac{F}{2\omega}\,t\sin(\omega t)$ 是
        一个解，并由第
        14 问中 $\gamma \to \omega$ 时的极限导出它：振幅
        线性地无限增长 --- 即共振灾变。
  \item （与傅里叶的联系）一般周期外力可分解为
        若干谐波（见傅里叶分析一章）；由线性性，
        稳态即各谐波响应之和。对于
        频率 $\omega = 3$ 的无阻尼振子，
        受方波驱动，该方波属于
        \cref{exo:b2:fourier:1}类型（谐波取
        全体奇整数），哪一次谐波会发生共振？用一句话说明
        工程师为何畏惧方波。
\end{enumerate}

\medskip
\noindent\textbf{第 IV 部分 --- 施图姆定理。} 考虑方程 $y''
+ q(t)\,y = 0$ 在区间 $I$ 上，$q$ 连续。（任一
方程 $y'' + ay' + by = 0$ 可借助
代换 $y = u\exp\bigl(-\frac12\int a\bigr)$ 化为此标准形式；
第 21 题展示了这一技巧的一个变形应用。）
\begin{enumerate}[resume]
  \item 对两解 $y_1, y_2$，证明朗斯基行列式
        $W = y_1y_2' - y_1'y_2$ 是\emph{常数}，
        且其为零当且仅当两解成比例；
        并证明非零解的零点都是\emph{简单、孤立}的。
  \item （施图姆分离定理）设 $y_1, y_2$ 为线性无关的解，
        $a < b$ 为 $y_1$ 的两个相邻零点。
        证明 $y_2$ 在
        $\intoo{a}{b}$ \emph{（计算常数 $W$ 在 $a$
        与 $b$ 处的值：此处 $W = y_1'y_2$，且 $y_1'(a)$
        与 $y_1'(b)$ 异号）}：线性无关解的
        零点彼此交错。
  \item （施图姆比较定理）设 $q_1 \leq q_2$ 在 $I$ 上，$y \neq
        0$ 满足 $y'' + q_1y = 0$，$z \neq 0$ 满足 $z'' + q_2z
        = 0$，$a < b$ 为 $y$ 的两个相邻零点。证明
        $z$ 在 $\intcc{a}{b}$ 中有零点 --- 若
        $q_1 < q_2$ 于 $\intoo ab$ 某处成立，则零点严格位于内部\emph{（若 $z \neq
        0$ 在 $\intoo ab$ 上，研究 $(yz' - y'z)' = (q_1 -
        q_2)yz$，对 $y, z$ 取定符号，并比较
        边界值）}。
  \item 由此推出\emph{间距界}：若 $0 < m^2 \leq
        q(t) \leq M^2$ 在 $I$ 上，则任意两个相邻零点
        $a < b$ （属于 $y'' + qy = 0$ 的非零解）满足
        下式：
        \[
        \frac{\pi}{M} \;\leq\; b - a \;\leq\; \frac{\pi}{m}
        \]
        \emph{（与 $u'' + M^2u = 0$ 及 $u'' + m^2u
        = 0$，其零点间距分别为 $\frac\pi M$ 与
        $\frac\pi m$）}。在调和振子上验证。
  \item 将 $ty'' + 2y' + ty = 0$
        （见\cref{exo:b2:diffeq:7}），令 $u = ty$ 化为 $u'' + u =
        0$，立即求出其解 $\frac{\sin t}t$，
        $\frac{\cos t}{t}$，并得出每个非零解的零点
        间距恰为
        $\pi$：施图姆的世界观 --- 零点由
        系数 $q$ 控制，无论有无显式公式。
\end{enumerate}

\medskip
\noindent\textbf{第 V 部分 --- 杜哈梅原理与有界性的
临界。}
\begin{enumerate}[resume]
  \item （振子的杜哈梅原理）证明：
        对连续的 $F$，方程 $x'' + \omega^2x =
        F(t)$ 满足 $x(0) = x'(0) = 0$ 的解为
        \[
        x(t) = \frac1\omega\int_0^t\sin\bigl(\omega(t -
        s)\bigr)F(s)\,\dd s ,
        \]
        并由此重新推导第 15 题中的共振解，
        取 $F(s) = F\cos(\omega s)$ 代入
        \emph{（积化和差）}。
  \item （$\zeta > 0$：有界输入、有界输出）证明
        当 $\zeta > 0$ 时，对\emph{任意}有界连续
        函数 $F$，阻尼方程的每个解均在
        $\intco{0}{\infty}$ 上有界\emph{（用常数变易法，并利用
        $\vertiii{\eu^{tA}} \leq
        C\eu^{-\alpha t}$ 的指数衰减性质，见\cref{exo:b2:diffeq:8}）}。
  \item （$\zeta = 0$）证明：无阻尼时，有界
        周期外力使所有解保持有界，
        \emph{除了}恰在共振时（$\gamma =
        \omega$，对比第 15 题与第 14 题）：阻尼的作用是
        把有界性的临界转化为一致
        稳定性。
  \item 综合。每项用一句话说明：（i）迹--行列式平面
        如何组织第 I--II 部分，以及
        强迫项（第 III 部分）将其引向何处；（ii）
        $\gamma_*$、$A(\gamma_*)$ 与
        因子 $\frac1{2\zeta}$ 的物理含义；（iii）施图姆定理所断言的
        正是显式公式无法给出的内容；（iv）
        本题中哪两个结论将在全书其余部分被悄然沿用
        （刘维尔常值朗斯基行列式；
        有界输入稳定性）。
\end{enumerate}
\end{problem}
